% !PRINTED 489
\weberpage{489}
If we subject the indices of the variables $u$, or, what comes to the same thing, the $u$ themselves to a permutation $\pi$, then the function $\psi$ may change, or it may remain unchanged. If it remains unchanged, we say that it \gesperrt{admits the permutation $\pi$}. If it admits all $\Pi(m)$ permutations that are possible at all, then it is symmetric and can be expressed by the symmetric fundamental functions (\S.\,44). The other extreme case is that in which the function takes nothing but distinct values for all $\Pi(m)$ arrangements. In general, certain permutations will leave the function $\psi$ unchanged, and others will change it. We now set up the theorem:

1. \gesperrt{The totality of all those permutations of $u, u_1, u_2 \ldots u_{m-1}$ which leave an integral function $\psi(u, u_1, u_2 \ldots u_{m-1})$ unchanged is a group of permutations.}

To prove it, assume that $\pi$ is one of the permutations which leaves $\psi$ unchanged. We express this by the following identity.
\[
\psi = \psi_\pi
\]
Now since the $u$ are independent variables, this identical equation must remain correct when the variables are subjected to any permutation $\pi'$ whatever. But when $\pi'$ is applied to $\psi_\pi$, the result is the same as when the composed permutation $\pi\pi'$ is applied to $\psi$. It follows therefore that
\[
\psi_{\pi'} = \psi_{\pi\pi'}
\]
and if now
\[
\psi = \psi_{\pi'}
\]
then it follows that
\[
\psi = \psi_{\pi\pi'},
\]
i.\,e. if the function $\psi$ remains unchanged under the permutations $\pi$ and $\pi'$, then it also remains unchanged under $\pi\pi'$, whereby, according to the definition of the group, the theorem we have stated is proved.

We have in this a means of forming groups of permutations, by taking any function of $m$ variables and seeking out all the permutations that leave such a function unchanged. That all groups of permutations can be formed in this way will appear later.
